\section{Background: how one release runs one file}
\label{sec:background}

This section describes the parts of the llama.cpp CPU backend that matter for the experiments.
Everything here refers to release b11696 (commit \texttt{ceefd16da}, 16 October 2026); file
names are relative to \texttt{ggml/src/ggml-cpu/}.

\paragraph{One release, several kernel sets.}
The prebuilt Windows CPU package of a llama.cpp release ships one \texttt{ggml-cpu-*.dll} per
x86 microarchitecture tier (\texttt{x64}, \texttt{sse42}, \texttt{sandybridge},
\texttt{ivybridge}, \texttt{haswell}, \texttt{alderlake}, \texttt{skylakex}, \texttt{icelake},
\texttt{zen4} and others). At start-up the loader scores each library against the CPU's feature
flags and loads the best one. The user does not choose and is not told, apart from a
\texttt{system\_info} log line. Two people who download the same release and the same GGUF file
therefore run different compiled kernels if their CPUs differ. Because the selection only
looks at which libraries are present, a single machine can run any tier its CPU supports by
leaving only that library in the folder. This is how the ISA experiments in this paper emulate
older or simpler hardware on one laptop.

\paragraph{Where the ISA enters the arithmetic.}
Quantized matrix products split into an integer part and a floating-point part. For a Q4\_0
weight block and a Q8\_0 activation block, the 32 products of 4-bit and 8-bit integers are
summed in integer arithmetic, which is exact on every instruction set. The block result is
then scaled by the two block scales and accumulated across the row in single precision.
That float accumulation is where tiers differ. In \texttt{arch/x86/quants.c} the AVX2 branch
of \texttt{ggml\_vec\_dot\_q4\_0\_q8\_0} accumulates with \texttt{\_mm256\_fmadd\_ps}, a fused
multiply-add with one rounding; the AVX branch uses \texttt{\_mm256\_add\_ps(\_mm256\_mul\_ps(\dots))},
which rounds twice; the SSSE3 and generic branches use other lane layouts and summation orders.
The same split appears in the Q4\_K, Q6\_K and F16 dot products and in the vector helpers used
by normalization and softmax (\texttt{simd-mappings.h} defines \texttt{GGML\_F32x8\_FMA} as
\texttt{\_mm256\_fmadd\_ps} only when \texttt{\_\_FMA\_\_} is defined).

\paragraph{Weight repacking.}
On CPUs that report AVX2, Q4\_0 and Q4\_K weight matrices are copied at load time into an
interleaved $8\times8$ layout (\texttt{repack.cpp},
\texttt{ggml\_repack\_get\_optimal\_repack\_type}); Q8\_0, Q6\_K and F16 are not repacked on
x86. The repacked matmul processes the activation rows of a micro-batch in groups of four with
a \texttt{gemm} kernel and handles the remaining $n \bmod 4$ rows one at a time with a
\texttt{gemv} kernel (\texttt{forward\_mul\_mat\_one\_chunk}). If the two kernels round
differently, a token's logits depend on its row index within its
micro-batch modulo four, which is a function of prompt length, micro-batch size, prompt-cache
hits and which other requests share the batch.

\paragraph{Prompt path versus generation path.}
For multi-row products llama.cpp calls tinyBLAS (\texttt{llamafile\_sgemm}) when the
shape qualifies: never for a single row (\texttt{n < 2}), and for Q8\_0 and Q4\_0 only when
$n \ge 8$ or $n = 4$; F16 additionally needs F16C. Otherwise it falls back to the per-row
\texttt{vec\_dot}. Attention's $QK^\top$ product reads the F16 key cache and switches path
with the number of query rows in the same way. The llama.cpp server documentation already warns
that logits are not guaranteed to be bit-identical between prompt processing and token
generation, and that the prompt cache can therefore make results nondeterministic.

\paragraph{Knobs a user treats as neutral.}
The following settings are documented as performance or memory options and do not change the
model: the thread count (\texttt{-t}), the logical and physical batch sizes (\texttt{-b},
\texttt{-ub}), flash attention (\texttt{-fa}), weight repacking (\texttt{--repack}), the number
of server slots (\texttt{-np}) and the prompt cache (\texttt{cache\_prompt}, on by default in the
server). In exact arithmetic every one of them computes the same function. I call two
configurations \emph{execution-equivalent} when they differ only in such settings or in the
selected ISA tier, and use the same release, the same GGUF file, the same prompt and greedy
decoding.
